\section{Direktif Data dan Alur Translasi}
\label{sec:direktif-translasi}

Di dalam berkas \file{.ASM} terdapat dua jenis pernyataan yang harus
dibedakan. \term{Instruksi} adalah perintah yang akan dijalankan
prosesor, misalnya \asm{MOV} atau \asm{ADD}. \term{Direktif} adalah
perintah kepada assembler, misalnya untuk menyediakan data atau
menentukan model memori. Direktif tidak diterjemahkan menjadi
instruksi yang dieksekusi.

Bagian ini diringkas dari \cite{gunawan2019}, Bagian 4,
hlm.~26--34, dan Bagian 8, hlm.~76--77, serta \cite{sto2001},
BAB V--VI, hlm.~21--26, dan BAB IX, hlm.~41--46.

\subsection{Direktif Data}

Direktif data menentukan isi segmen data. Bentuk yang paling sering
dipakai dirangkum pada Tabel~\ref{tab:direktif-data}.

\begin{table}[H]
    \centering
    \caption{Direktif data yang sering dipakai pada program 8086.}
    \label{tab:direktif-data}
    \small
    \begin{tabular}{@{}lp{3.8cm}p{7.0cm}@{}}
        \toprule
        \textbf{Direktif} & \textbf{Ukuran} & \textbf{Kegunaan} \\
        \midrule
        \dir{DB} & 1 byte & Mendefinisikan byte, karakter, atau
        string \\
        \addlinespace
        \dir{DW} & 2 byte & Mendefinisikan satu kata (\eng{word}) \\
        \addlinespace
        \dir{DD} & 4 byte & Mendefinisikan bilangan 32 bit \\
        \addlinespace
        \dir{EQU} & --- & Memberi nama pada konstanta \\
        \addlinespace
        \dir{DUP} & sesuai data & Mengulang pola data beberapa kali \\
        \addlinespace
        \code{?} & sesuai direktif & Menyediakan ruang tanpa nilai
        awal \\
        \addlinespace
        \dir{ORG} & --- & Menetapkan offset awal program atau data \\
        \bottomrule
    \end{tabular}
\end{table}

Contoh pemakaian:

\begin{lstlisting}[language=MASM, numbers=none]
pesan   DB  'Halo', 0Dh, 0Ah, '$'   ; string dan kode kendali
nilai   DW  1234h                    ; satu kata
besar   DD  12345678h                ; empat byte
PANJANG EQU 20                       ; konstanta
bintang DB  4 DUP('*')               ; empat karakter bintang
buffer  DB  8 DUP(?)                 ; ruang kosong delapan byte
\end{lstlisting}

String pada \dir{DB} disimpan sebagai byte-byte berurutan. Karakter
\code{0Dh} dan \code{0Ah} dipakai untuk pindah baris, sedangkan
\code{'\$'} menandai akhir string untuk layanan DOS fungsi
\hex{09}. Konstanta \dir{EQU} tidak menempati memori; nilainya
disisipkan assembler ketika nama tersebut dipakai.

\subsection{Model Memori dan Direktif Struktur}

Direktif \dir{.MODEL} menggantikan penulisan \dir{SEGMENT} dan
\dir{ASSUME} yang panjang pada program sederhana. Model \code{SMALL}
menyediakan satu segmen kode dan satu segmen data. Model \code{TINY}
menggabungkan kode, data, dan \eng{stack} dalam satu segmen untuk
program \file{.COM}.

Direktif \dir{ORG 100h} pada program \file{.COM} menetapkan bahwa
program dimuat mulai offset \hex{100}. Pada program \file{.EXE},
direktif ini tidak dipakai karena linker yang menentukan alamat awal
berdasarkan informasi program.

\begin{infobox}
\vspace{1.1em}
\important{Direktif bukan instruksi.} \dir{DB}, \dir{DW}, \dir{EQU},
dan \dir{.MODEL} tidak menjadi instruksi mesin. Direktif hanya
mengatur bagaimana assembler menyusun program dan data.
\end{infobox}

\subsection{Alur Translasi dan Berkas Hasil}

Alur translasi program dapat diringkas sebagai berikut.

\begin{enumerate}
    \item \file{program.asm} ditulis dalam editor teks.
    \item TASM menghasilkan \file{program.obj}.
    \item TLINK menghasilkan \file{program.exe} atau, dengan pilihan
          \code{/t}, \file{program.com}.
    \item DOS memuat berkas hasil dan menjalankannya.
\end{enumerate}

Berkas \file{.OBJ} belum dapat dijalankan karena masih memuat
rujukan simbol dan alamat yang belum final. TLINK menyelesaikan
rujukan tersebut dan menetapkan alamat segmen. Berkas \file{.EXE}
memuat informasi tambahan tentang segmen dan titik masuk, sedangkan
berkas \file{.COM} dimuat langsung pada satu segmen.

\subsection{Membaca Pesan Kesalahan Assembler}

TASM melaporkan kesalahan beserta nomor barisnya. Pesan yang sering
muncul pada tahap awal dirangkum pada Tabel~\ref{tab:error-assembler}.

\begin{table}[H]
    \centering
    \caption{Pesan kesalahan assembler yang sering muncul.}
    \label{tab:error-assembler}
    \small
    \begin{tabular}{@{}p{4.2cm}p{4.4cm}p{5.0cm}@{}}
        \toprule
        \textbf{Pesan} & \textbf{Penyebab umum} &
        \textbf{Cara memperbaiki} \\
        \midrule
        \file{Undefined symbol} & Label atau variabel belum
        didefinisikan & Periksa ejaan nama dan pastikan deklarasinya
        ada \\
        \addlinespace
        \file{Syntax error} & Salah tulis instruksi atau direktif &
        Periksa penulisan mnemonic, koma, dan tanda kutip \\
        \addlinespace
        \file{Operand types do not match} & Ukuran operand tidak
        sama & Samakan ukuran register dan data, atau gunakan
        \lstasm|BYTE PTR| / \lstasm|WORD PTR| \\
        \addlinespace
        \file{Improper operand type} & Operasi memori ke memori yang
        tidak diizinkan & Pindahkan salah satu operand ke register \\
        \bottomrule
    \end{tabular}
\end{table}

\begin{peringatanbox}
\vspace{1.1em}
Kesalahan yang tidak dilaporkan assembler tetapi membuat program
gagal berjalan biasanya berasal dari logika: salah memilih register,
lupa menginisialisasi \reg{DS}, atau salah menghitung offset. Karena
itu, pesan kesalahan assembler hanya menangani kesalahan penulisan;
kesalahan logika harus dicari dengan \eng{debugger}.
\end{peringatanbox}
